\section{FUSIONATTN: Multi-Paradigm Hybrid Sparse Attention with Block Gating and Short-Sequence Fast Path}

\textbf{Authors:}
\begin{itemize}
\item {{AUTHOR\_1\_NAME}}  -- {{AUTHOR\_1\_EMAIL}}
\end{itemize}
\begin{itemize}
\item {{AUTHOR\_2\_NAME}}  -- {{AUTHOR\_2\_EMAIL}}
\end{itemize}

\textbf{Affiliations:}
\begin{itemize}
\item   {{UNIVERSITY\_1}}, {{DEPARTMENT\_1}}
\end{itemize}
\begin{itemize}
\item   {{UNIVERSITY\_2}}, {{DEPARTMENT\_2}}
\end{itemize}

\textbf{Corresponding author:} {{CORRESPONDING\_AUTHOR}}
\textbf{Preprint date:} {{DATE}}
\textbf{Version:} v{{VERSION}}
\textbf{Code:} https://github.com/{{ORG}}/SVD\_Project/tree/main/framework/attentions/fusion\_attn.py
\textbf{License:} {{LICENSE}}



\subsection{Abstract}

Long-context attention in transformers is quadratic in sequence length, 
$\mathcal{O}(L^2 d)$, which makes inference prohibitively expensive at 
$L \geq 4096$. Existing sparse methods -- Longformer, BigBird, MoBA -- reduce 
this to $\mathcal{O}(L \cdot W \cdot d)$ or $\mathcal{O}(K \cdot L \cdot d)$ 
but often apply the same sparsity pattern regardless of sequence length. 
We introduce \textbf{FUSIONATTN}, a \textbf{multi-paradigm hybrid attention} that 
dispatches between two paths based on sequence length:

1. \textbf{Short-sequence path ($L \leq 128$):} Standard softmax attention, no 
   sparsity overhead.
2. \textbf{Long-sequence path ($L > 128$):} Block-sparse attention with a 
   \textbf{learnable gating MLP} that scores query-key block pairs, selects the 
   top-$K$ blocks per query block, and computes attention only within the 
   selected blocks (\textbf{MoBA-style}).

The gate is a 2-layer MLP with GELU activation that operates on 
mean-pooled block representations; it is trained jointly with LoRA 
adapters. Base Q/K/V/O projections are frozen and copied exactly from the 
pretrained model, ensuring step-0 preservation.

\textbf{Important disclosure.} FUSIONATTN is \textbf{not} a compression method. It 
\textbf{increases} parameter count via LoRA adapters and the gate MLP. Its value 
lies in \textbf{long-context compute efficiency}, not size reduction. We 
additionally note that at $L \leq 128$ (our test sequence length), FUSIONATTN 
takes the \textbf{short path} and produces output identical to the pretrained 
model with no sparsity benefit.

We evaluate on DistilBERT, RoBERTa-base, and GPT-2 on SST-2 with 
{{RECOVERY\_STEPS}} LoRA recovery steps. FUSIONATTN achieves 
\textbf{{{D\_ACC}}\%} on DistilBERT, \textbf{{{R\_ACC}}\%} on RoBERTa, and \textbf{{{G\_ACC}}\%} 
on GPT-2, matching the SDPA baseline within noise.

\textbf{Keywords:} hybrid sparse attention, block gating, MoBA, short-sequence 
fast path, learnable sparsity, LoRA



\subsection{Lay Summary}

Standard attention compares every token to every other token, which becomes 
slow for long sequences. FUSIONATTN takes a hybrid approach: for short 
sequences (under 128 tokens), it uses the standard fast attention. For 
longer sequences, it groups tokens into blocks and uses a small neural 
network to decide which blocks should "talk to" each other. Only the 
selected blocks compute attention. This saves computation for long 
sequences while keeping short-sequence behavior unchanged. The trade-off: 
it adds some parameters to the model, and it's not faster on short 
sequences.



\subsection{Table of Contents}

1. [Introduction](01\_intro.md)
2. [Related Work](02\_related.md)
3. [Background](03\_background.md)
4. [Method](04\_method.md)
5. [Architecture Integration](05\_architecture.md)
6. [Experimental Setup](06\_setup.md)
7. [Results](07\_results.md)
8. [Analysis](08\_analysis.md)
9. [Discussion](09\_discussion.md)
10. [Limitations](10\_limitations.md)
11. [Conclusion](11\_conclusion.md)
12. [Reproducibility](12\_reproducibility.md)
13. [References](13\_references.md)

\textbf{Appendices:}
\begin{itemize}
\item [A. Mathematical Derivations](A\_math.md)
\end{itemize}
\begin{itemize}
\item [B. Implementation Details](B\_implementation.md)
\end{itemize}
\begin{itemize}
\item [C. Full Experimental Data](C\_data.md)
\end{itemize}
\begin{itemize}
\item [D. Negative Results](D\_negative\_results.md)
\end{itemize}
\begin{itemize}
\item [E. Ethical Considerations](E\_ethics.md)
\end{itemize}